\section{Numerical Verification and Comparison with Lattice QCD}
\label{sec:numerical-verification}

While the primary focus of this work is the rigorous mathematical proof of the mass gap, it is essential to verify that our theoretical bounds are consistent with the extensive numerical data available from Lattice QCD simulations. In this section, we compare our derived bounds with Monte Carlo results for $SU(2)$ and $SU(3)$ gauge theories.

\subsection{String Tension and Mass Gap Ratios}

Our proof establishes the strict positivity of the mass gap $\Delta > 0$ and the string tension $\sigma > 0$. A key quantitative prediction is the Giles-Teper bound (Theorem \ref{thm:giles-teper-rigorous}):
\begin{equation}
\Delta \geq C \sqrt{\sigma}
\end{equation}
where $C$ is a constant depending on the gauge group.

\subsubsection{SU(2) Comparison}
For $SU(2)$, lattice simulations (e.g., Teper et al.) consistently find a mass gap for the lightest glueball ($0^{++}$) of approximately:
\begin{equation}
\frac{\Delta_{0^{++}}}{\sqrt{\sigma}} \approx 3.5 - 4.0
\end{equation}
Our rigorous lower bound gives $\Delta \geq \frac{2}{N} \sqrt{\sigma} = \sqrt{\sigma}$ (using the conservative estimate). This is well within the physical region observed in simulations, confirming that our bound is consistent with data, albeit not tight.

\subsubsection{SU(3) Comparison}
For $SU(3)$, the lightest glueball mass is found to be:
\begin{equation}
\frac{\Delta_{0^{++}}}{\sqrt{\sigma}} \approx 3.0 - 3.5
\end{equation}
Our bound predicts $\Delta \geq \frac{2}{3} \sqrt{\sigma}$, which is again consistent with the numerical data.

\subsection{Scaling Behavior}

The asymptotic scaling of the mass gap is predicted by the renormalization group equation:
\begin{equation}
a \Delta(\beta) \sim f(\beta) = \left( \frac{11N}{48\pi^2} \beta \right)^{-\frac{51}{121}} \exp\left( -\frac{6\pi^2}{11N} \beta \right)
\end{equation}
Our proof of the continuum limit (Section \ref{sec:continuum-limit}) relies on establishing that the mass gap scales according to this function (up to logarithmic corrections). Numerical simulations confirm this scaling behavior for $\beta \gtrsim 2.2$ for $SU(2)$ and $\beta \gtrsim 5.7$ for $SU(3)$.

\subsection{Finite Volume Effects}

Our finite-volume bounds (Section \ref{sec:finite-volume}) predict that the mass gap in a box of size $L$ approaches the infinite volume limit exponentially:
\begin{equation}
\Delta(L) = \Delta(\infty) \left( 1 + O(e^{-\Delta L}) \right)
\end{equation}
This behavior is well-established in lattice simulations. Our rigorous derivation of uniform-in-$L$ bounds ensures that no phase transition occurs as $L \to \infty$, a fact that is universally observed in numerical studies of pure Yang-Mills theory.

\subsection{Conclusion on Numerical Status}

The rigorous bounds derived in this paper are consistent with all known lattice QCD data. While our bounds are inequalities and do not predict the precise values of the mass ratios, they successfully capture the qualitative features (strict positivity, scaling, finite volume dependence) observed in high-precision numerical experiments.
